\clearpage
\section*{Appendix R\quad Computational Receipts}\label{app:receipts}
\addcontentsline{toc}{section}{Appendix R: Computational Receipts}
\small\noindent Each computation marked \rcptmarker\ in the text is verified by a runnable receipt in the \texttt{receipts/} directory that ships with this work; status \textsf{OK} = verified (traced, run, output matches the claim). This appendix is generated from the receipt index, not maintained by hand.\par\medskip
\begingroup\renewcommand{\arraystretch}{1.25}
\begin{description}
\item[\label{rcpt:X5_monodromy_group}\texttt{X5\_monodromy\_group}]\hfill\textsf{[OK]}\\
\textit{`rem:monodromy-group` --- **the generation is verified, and the base has three punctures**} \ (P5 `rem:monodromy-group` --- **the generation is verified, and the base has three punctures**). 
\emph{Computes:} run r1871, rc=0
 \emph{Bound:} **bound, and it records a caught error: a first pass computed \$c\_\textbackslash\{\}infty\$ based at \$m=60\$ rather than the common base point, and the relation appeared to fail. Permutations compose only when all loops share a base point**
\ \emph{Run:} \texttt{python3 receipts/storyboard\_receipts/X5\_monodromy\_group.py}
\item[\label{rcpt:X1rr_sky_angle_is_the_galois_closure}\texttt{X1rr\_sky\_angle\_is\_the\_galois\_closure}]\hfill\textsf{[OK]}\\
\textit{`rem:galois-closure` --- **the dial is the splitting cover**} \ (P5 `rem:galois-closure` --- **the dial is the splitting cover**). 
\emph{Computes:} run r1881, rc=0
 \emph{Bound:} ***supersedes `X1r\_cover\_tower`, which mis-identified both the deck group and the covering involution***
\ \emph{Run:} \texttt{python3 receipts/storyboard\_receipts/X1rr\_sky\_angle\_is\_the\_galois\_closure.py}
\item[\label{rcpt:C1_conformal_vs_isometric}\texttt{C1\_conformal\_vs\_isometric}]\hfill\textsf{[OK]}\\
\textit{alpha-action --- **the homothety IS the conformal dilation**} \ (P5 \S{}alpha-action --- **the homothety IS the conformal dilation**). 
\emph{Computes:} run r1853, rc=0
 \emph{Bound:} **bound: it computes the LINEAR conformal maps. The full conformal group of the null cone includes special conformal transformations not tested here**
\ \emph{Run:} \texttt{python3 receipts/storyboard\_receipts/C1\_conformal\_vs\_isometric.py}
\item[\label{rcpt:Q4_equianharmonic_vantages}\texttt{Q4\_equianharmonic\_vantages}]\hfill\textsf{[OK]}\\
\textit{`rem:equianharmonic` --- **the groupoid's projective invariant**} \ (P5 `rem:equianharmonic` --- **the groupoid's projective invariant**). 
\emph{Computes:} run r1850, rc=0
 \emph{Bound:} **bound, stated in the script: three points \$120\textasciicircum{}\textbackslash\{\}circ\$ apart with their centre are equianharmonic AUTOMATICALLY --- the content is that CR's vantages ARE so spaced, which P3 forces. It names a consequence, not a new fact**
\ \emph{Run:} \texttt{python3 receipts/storyboard\_receipts/Q4\_equianharmonic\_vantages.py}
\item[\label{rcpt:negation_outer_A2}\texttt{negation\_outer\_A2}]\hfill\textsf{[OK]}\\
\textit{rem:orientation} \ (R negates the roots (r->-r) conjugates the rep 3<->3bar ONLY because -1 not in W(A2); -1 in W for A1xA1/B2/G2 (inner)). 
\emph{Computes:} Weyl groups built from simple reflections; -1 membership across 4 rank-2 systems (others = controls)
 \emph{Bound:} shares the -1 not-in-W(A2) core with P14 two\_factors\_direct (distinct claim, kept self-contained)
\ \emph{Run:} \texttt{python3 receipts/P05\_groupoid/negation\_outer\_A2.py}
\item[\label{rcpt:P05_deck_group_S3}\texttt{P05\_deck\_group\_S3}]\hfill\textsf{[OK]}\\
\textit{sec:deck/sec:classification} \ (deck group of the 3-sheeted cover = S3 (Nariai monodromies = transpositions, generate S3 = Galois group); with orientation Z2 = D6). 
\emph{Computes:} analytic continuation of cubic roots around each Nariai branch (loop from common base) = distinct transpositions -> S3; discriminant not a square -> Galois S3; S3xZ2=D6 order 12; control: generic loop trivial
 \emph{Bound:} ---
\ \emph{Run:} \texttt{python3 receipts/P05\_groupoid/P05\_deck\_group\_S3.py}
\item[\label{rcpt:P05_dihedral_generators}\texttt{P05\_dihedral\_generators}]\hfill\textsf{[OK]}\\
\textit{prop:relations/completeness} \ (within-single-geometry morphism group = dihedral sigma,tau on the sky-angle circle: sigma\textasciicircum{}2=tau\textasciicircum{}3=(sigma tau)\textasciicircum{}2=id, <sigma,tau> order 6 = D3\textasciitilde{}S3, all 6 preserve sin 3w). 
\emph{Computes:} builds generators, checks relations + order 6 + sin3w-invariance
\ \emph{Run:} \texttt{python3 receipts/P05\_groupoid/P05\_dihedral\_generators.py}
\item[\label{rcpt:COVER_dial_loop_is_sigma_squared}\texttt{COVER\_dial\_loop\_is\_sigma\_squared}]\hfill\textsf{[OK]}\\
\textit{rem:perroot} \ (How does a loop in the DIAL project to the MASS plane? 2M = (2/3sqrt3) a sin 3w. Near a branch point sin3w = +-1 the map is critical, so the projection winds twice.). 
\emph{Computes:} runs clean; registered r2376 (c54) from its own docstring and a live run
 \emph{Bound:} description is the receipt's own wording, condensed; not independently re-derived here
\ \emph{Run:} \texttt{python3 receipts/P05\_groupoid/COVER\_dial\_loop\_is\_sigma\_squared.py}
\item[\label{rcpt:HOLONOMY_complex_loop}\texttt{HOLONOMY\_complex\_loop}]\hfill\textsf{[OK]}\\
\textit{rem:perroot} \ (Transport the RESIDUE FORM around a Nariai point in the COMPLEX w-plane. The roots return (splitting field, no monodromy) -- but the FORM's transport can accumulate. A = -(1/2) G\textasciicircum{}\{-1\} dG with G\_ii = 1/f'(r\_i); f' has a simple zero at the double root -> simple pole in A.). 
\emph{Computes:} runs clean; registered r2376 (c54) from its own docstring and a live run
 \emph{Bound:} description is the receipt's own wording, condensed; not independently re-derived here
\ \emph{Run:} \texttt{python3 receipts/P05\_groupoid/HOLONOMY\_complex\_loop.py}
\item[\label{rcpt:K1_action_groupoid}\texttt{K1\_action\_groupoid}]\hfill\textsf{[OK]}\\
\textit{sec:nariai-algebroid} \ (K1 --- WHAT GROUPOID DOES P12's ALGEBROID INTEGRATE TO? CONFIRMATION 4, from P12 at source: base = the SPACE OF CUTS C of dS\_5 = SO(5,1)/SO(4,1) section = the cosmic clock (P10) structure function = the inverse spatial metric and the object P12 names in its own notation: 'what the continuous so(5,1) \textbar{}x C realizes' connection = the leak of [m,m] into m, 'vanishing exactly at the two symmetric strata' P5's groupoid, by c...). 
\emph{Computes:} runs clean; registered r2376 (c54) from its own docstring and a live run
 \emph{Bound:} description is the receipt's own wording, condensed; not independently re-derived here
\ \emph{Run:} \texttt{python3 receipts/P05\_groupoid/K1\_action\_groupoid.py}
\item[\label{rcpt:K45_span_and_nariai}\texttt{K45\_span\_and\_nariai}]\hfill\textsf{[OK]}\\
\textit{sec:nariai-algebroid} \ (K4 --- THE THREE LEVELS AS A DIAGRAM. K5 --- IS THE NARIAI COINCIDENCE FORCED? CONFIRMATION 1/3 from section 1-LEVELS, at source: L1 the foliation stacking rate (the existent's own geometric expansion) \ensuremath{\cdot} L2 the leaf's self-gravitating dynamics \ensuremath{\cdot} L3 the E=1 projection (the shadow).). 
\emph{Computes:} runs clean; registered r2376 (c54) from its own docstring and a live run
 \emph{Bound:} description is the receipt's own wording, condensed; not independently re-derived here
\ \emph{Run:} \texttt{python3 receipts/P05\_groupoid/K45\_span\_and\_nariai.py}
\item[\label{rcpt:V5_homothety_charge}\texttt{V5\_homothety\_charge}]\hfill\textsf{[OK]}\\
\textit{sec:alpha-action} \ (V5 --- THE DILATION AS A SYMMETRY WITH NO CONSERVED CHARGE. CONFIRMATION 4 --- the object: C1 established that the group preserving the substrate's ABSOLUTE is O(5,1) x R+, one generator larger than the isometry group O(5,1), and that the extra generator is the DILATION carrying alpha -> lambda.alpha. A dilation is a HOMOTHETY: L\_xi g = 2c g with c constant.). 
\emph{Computes:} runs clean; registered r2376 (c54) from its own docstring and a live run
 \emph{Bound:} description is the receipt's own wording, condensed; not independently re-derived here
\ \emph{Run:} \texttt{python3 receipts/P05\_groupoid/V5\_homothety\_charge.py}
\item[\label{rcpt:P03_acceleration_is_slice_curvature}\texttt{P03\_acceleration\_is\_slice\_curvature}]\hfill\textsf{[OK]}\\
\textit{sec:curvature \ensuremath{\cdot} rem:twocritical \ensuremath{\cdot} the bend} \ (on ANY member of the energy family (E cancels: (dr/dtau)\textasciicircum{}2 = E\textasciicircum{}2 - f, and E is a constant of the motion) d\textasciicircum{}2r/dtau\textasciicircum{}2 = -f'/2 = r*K\_G identically, so THE SIGN OF THE SLICE'S GAUSSIAN CURVATURE IS THE SIGN OF THE COMOVING ACCELERATION -- and the flat locus is the deceleration-to-acceleration turnover, which on Nariai is the seam and in standard terms is rho\_m = 2 rho\_Lambda (csch\textasciicircum{}2 w = 2 <=> sinh w = 1/sqrt2, the seam's phase). A recent cosmological epoch: 7.06 Gyr at H0=73, 7.65 at 67.4, preceding matter-Lambda equality [borrowed from P3 \ensuremath{\cdot} P7 \ensuremath{\cdot} P8]). 
\emph{Computes:} the identity asserted; the K\_G zero solved and matched to (M alpha\textasciicircum{}2)\textasciicircum{}(1/3); csch\textasciicircum{}2 w = 2 <=> sinh w = 1/sqrt2 proved by algebra after simplify() failed on the closed forms; epochs at both H0
 \emph{Bound:} bound: identities on the Nariai E=1 worldline; no causal claim between seam and turnover -- they are one locus; dating inherits H0
\ \emph{Run:} \texttt{python3 receipts/P03\_SdS\_slicing/P03\_acceleration\_is\_slice\_curvature.py}
\item[\label{rcpt:C5_sct_already_linear}\texttt{C5\_sct\_already\_linear}]\hfill\textsf{[OK]}\\
\textit{rem:reality-lines} \ (TWO CLOSURES. (A) X3r's open question: sin is real on the real axis AND on every line Re z = pi/2 + k.pi. The crossings are at z = pi/2 + k.pi. Are the odd ones the corpus's BACK seam as the even ones are its FRONT? (B) The conformal queue's last entry: 'C1 computed the LINEAR conformal maps only; the SPECIAL CONFORMAL TRANSFORMATIONS are untested -- the one place the conformal field could still bite.' [borrowed from P3 (+P5)]). 
\emph{Computes:} runs clean; registered r2376 (c54) from its own docstring and a live run
 \emph{Bound:} description is the receipt's own wording, condensed; not independently re-derived here
\ \emph{Run:} \texttt{python3 receipts/P03\_SdS\_slicing/C5\_sct\_already\_linear.py}
\item[\label{rcpt:K23_functor_and_quotient}\texttt{K23\_functor\_and\_quotient}]\hfill\textsf{[OK]}\\
\textit{rem:universal} \ (K2 + K3 --- IS Psi A FUNCTOR, AND IS 'THE SINGLE INVARIANT IS THE GEOMETRY' A UNIVERSAL PROPERTY? CONFIRMATION 4 --- the objects: SOURCE: K1's action groupoid SO(5,1) \textbar{}x C ==> C . Objects = cuts. Arrows c -> g.c . Psi (P8): 'carries a cut of the substrate to a geometry'. TARGET: geometries, with isometries between them. [borrowed from P5 (+P8)]). 
\emph{Computes:} runs clean; registered r2376 (c54) from its own docstring and a live run
 \emph{Bound:} description is the receipt's own wording, condensed; not independently re-derived here
\ \emph{Run:} \texttt{python3 receipts/P05\_groupoid/K23\_functor\_and\_quotient.py}
\end{description}\endgroup
